\section{Methodology}
\label{sec-ch2:methodology}
This work seeks to investigate smart building research to understand what researchers envision and design for. The goal is not to provide an exhaustive summary of existing research findings but to focus on identifying the human experiences that research in this domain aims to improve or introduce. In addition, we identified \emph{how} the field aims to shape those human experiences through design mechanisms and technological interventions. This was accomplished by conducting a scoping review based on the five-stage framework outlined by \citet{arksey2005scoping}. We chose a scoping review as this method suits broad research questions that map key characteristics rather than a targeted study of a specific effect \cite{munn2018systematic, rogers2024umbrella}. The review was planned with pre-defined research questions and a search strategy, rather than following a formal protocol, allowing the specifics of the analysis to be shaped by the first author's growing familiarity with the literature by reading the selected corpus papers. The process is documented in the rest of this section. 
To ensure comprehensive reporting, this paper broadly follows the Preferred Reporting Items for Systematic Reviews and Meta-Analyses extension for Scoping Reviews (PRISMA-ScR) checklist where applicable\sidenote{As conventions in HCI do not always accommodate the PRISMA-ScR items, we provide a table in the supplementary materials to detail the manner by which we adhere to the reporting checklist.}. The flow diagram in \autoref{fig:flowdiagram} provides an overview of the search and screening stages to reach the final review corpus. 


% \pgfdeclarelayer{background}
% \pgfsetlayers{background,main}

% \definecolor{green-include}{RGB}{102,194,164}

% \begin{mainfigure}[tbp]
% {
% \resizebox{\linewidth}{!}{
% \small
% \begin{tikzpicture}
%     \tikzstyle{block}=[
%         draw,
%         thick,
%         text width=2.5cm,
%         minimum height=1.2cm,
%         align=center
%     ]

%     \tikzstyle{header}=[rotate=90, align=center]
%     \tikzstyle{line}=[-latex, draw, very thick]

%     \node[block, dotted] (pilot) {pilot searches};

%     \node[block, right = 2em of pilot, minimum height=1.2cm, text height=1cm, minimum width=4.2cm, text width=3cm] (search) {identification};
%     \node[draw, align=center,  minimum width=1.2cm] (acm) at ([xshift=-1.3cm, yshift=.12cm]search){\small ACM DL\\n=3653};
%     \node[draw, align=center,  minimum width=1.2cm] (scopus) at ([yshift=.12cm]search){\small Scopus\\n=1873};
%     \node[draw, align=center,  minimum width=1.2cm] (other) at ([xshift=1.3cm, yshift=.12cm]search){\scriptsize \faSnowflake \\n=22};

%     \node[block, right=4em of search] (tascreen) {title and abstract screening};
    
%     \node[below=.8em of tascreen] (ie1) {\emph{criteria: building$+$people$+$technology}};

%     \node[block, right=4em of tascreen] (dedup) {deduplication};
    
%     \node[block, right=4em of dedup] (ftscreen) {full-text screening};

%     \node[below=.8em of ftscreen] (ie2) {\emph{criteria: see \autoref{sec:iecrit}}};

%     \node[block, below = 1.4em of dedup] (snowball) {snowballing};
    
%     \node[block, right=4em of ftscreen] (corpus) {final corpus: \textbf{N=192}};
%     % Draw edges

%     \path [line, dotted] (pilot) -- (search);

%     \path [line] (search) -- node [midway,above,align=center] {n=5548}(tascreen);

%     \path [line] (tascreen) -- node [midway, above, align=center] {n=368} (dedup);
    
%     \path [line] (dedup) -- node [midway, above, align=center] {n=338} (ftscreen);

    
%     \path [line] (ftscreen) -- node [midway, above, align=center] {n=192} (corpus);

%     \path [line, dotted] (ie2.south) |- ([xshift=2em]snowball.east) -- (snowball.east);
%     \path [line, dotted] (snowball.west) |- ([xshift=-5.75cm]snowball.west) -- (other.south);
    
%     %\path [line] (screening-0) |- ([yshift=.35cm]screening-2.north) -- (screening-2.north);
%     %\path [line] (identification-1.east) -- (identification-2.west);
%     %\path [line] (screening-0) -- (screening-1);
%     %\path [line] (screening-0) -- (screening-excl);
%     %\path [line] (screening-1) -- (screening-2);
%     %\path [line] (screening-1) -- (eligibility-1);
%     %\path [line] (eligibility-1) -- (eligibility-2);
%     %\path [line] (screening-1) |- ([yshift=.35cm]eligibility-2.north) -- (eligibility-2);
%     %\path [line] (eligibility-1) -- (included-0);
%     %\path [line] (identification-2) |- ([yshift=.3cm]screening-0.north) -- (screening-0.north);
    
    
%     %\path [line, draw=gray] (identification-3) -- (screening-3);
%     %\path [line, draw=gray] (screening-3) -- (eligibility-other);
%     %\path [line, draw=gray] (eligibility-other.south) |- ([yshift=.35cm, xshift=4cm]included-prefinal.north) -- ([yshift=.35cm]included-prefinal.north) -- (included-prefinal.north);
%     %\path [line] (included-0) -- (included-prefinal);
    
    
%     %\path [line] (included-prefinal) -- (included-final);
    
% \end{tikzpicture}
% }}
% \caption{This  flowchart demonstrates the scoping review process and the number of papers considered in each stage.}
% \label{fig:flowdiagram}
% \end{mainfigure}

% %\footnotetext{Based on the 2020 template: \url{https://www.prisma-statement.org/prisma-2020-flow-diagram}, last accessed on May 15, 2024}


\pgfdeclarelayer{background}
\pgfsetlayers{background,main}

\definecolor{green-include}{RGB}{102,194,164}

\begin{mainfigure}[tbp]
\centering

\resizebox{\linewidth}{!}{%
\begin{tikzpicture}[
    font=\footnotesize,
    >=Latex,
    node distance=2em and 4em,
    block/.style={
        draw,
        thick,
        align=center,
        minimum height=1.1cm,
        text width=2.4cm,
        inner sep=4pt
    },
    smallblock/.style={
        draw,
        align=center,
        minimum width=1.35cm,
        minimum height=.75cm,
        inner sep=2pt
    },
    line/.style={
        -{Latex[length=2mm]},
        draw,
        thick
    },
    auxline/.style={
        -{Latex[length=2mm]},
        draw,
        thick,
        dotted
    }
]

% Main nodes
\node[block, dotted, text width=2.2cm] (pilot) {pilot searches};

% \node[block, right=2em of pilot, minimum width=4.3cm, text width=3.2cm] (search) {identification};

% \node[smallblock] (acm) at ([xshift=-1.3cm,yshift=.10cm]search.center) {\small ACM DL\\n=3653};
% \node[smallblock] (scopus) at ([yshift=.10cm]search.center) {\small Scopus\\n=1873};
% \node[smallblock] (other) at ([xshift=1.3cm,yshift=.10cm]search.center) {\scriptsize \faSnowflake\\n=22};

% identification container
\node[
  draw,
  thick,
  minimum width=5.8cm,
  minimum height=1.35cm
] (search) [right=2em of pilot] {};

% label inside the outer box
\node[font=\small] at ([yshift=-0.34cm]search.center) {identification};

% three small boxes inside the identification box
\node[
  draw,
  align=center,
  minimum width=1.35cm,
  minimum height=0.58cm,
  font=\scriptsize,
  anchor=west
] (acm) at ([xshift=0.18cm,yshift=0.18cm]search.west)
{ACM DL\\n=3653};

\node[
  draw,
  align=center,
  minimum width=1.35cm,
  minimum height=0.58cm,
  font=\scriptsize,
  anchor=west
] (scopus) at ([xshift=1.78cm,yshift=0.18cm]search.west)
{Scopus\\n=1873};

\node[
  draw,
  align=center,
  minimum width=1.10cm,
  minimum height=0.58cm,
  font=\scriptsize,
  anchor=west
] (other) at ([xshift=3.38cm,yshift=0.18cm]search.west)
{\faSnowflake\\n=22};


\node[block, right=4em of search] (tascreen) {title and\\abstract screening};

\node[below=.8em of tascreen] (ie1) {\emph{criteria: building$+$people$+$technology}};

\node[block, right=4em of tascreen] (dedup) {deduplication};

\node[block, right=4em of dedup] (ftscreen) {full-text\\screening};

\node[below=.8em of ftscreen] (ie2) {\emph{criteria: see \autoref{sec:iecrit}}};

\node[block, below=1.4em of dedup] (snowball) {snowballing};

\node[block, right=4em of ftscreen] (corpus) {final corpus:\\\textbf{N=192}};

% Arrows
\draw[auxline] (pilot) -- (search);

\draw[line] (search) -- node[midway, above] {n=5548} (tascreen);
\draw[line] (tascreen) -- node[midway, above] {n=368} (dedup);
\draw[line] (dedup) -- node[midway, above] {n=338} (ftscreen);
\draw[line] (ftscreen) -- node[midway, above] {n=192} (corpus);

\draw[auxline] (ie2.south) |- ([xshift=2em]snowball.east) -- (snowball.east);
\draw[auxline] (snowball.west) |- ([xshift=-5.75cm]snowball.west) -- (other.south);

\end{tikzpicture}%
}

\caption{This flowchart demonstrates the scoping review process and the number of papers considered at each stage.}
\label{fig:flowdiagram}
\end{mainfigure}

\subsection{Framework Stage 1: Identifying Research Questions}
During the first stage in our scoping review work, key research questions were identified. The first question, established from the outset of the project, focuses on the different targeted human experiences in smart buildings:

\begin{itemize}
    \item \textbf{\textit{RQ1.}} What are the different types of \emph{human experiences} that research in the broad domain of smart buildings aims to improve or introduce?
\end{itemize}

The other two research questions underwent iterations but focused on the technologies and mechanisms that the field adopts from interaction design, architecture, and other areas to improve human experiences in smart buildings:  

\begin{itemize}
    \item \textbf{\textit{RQ2.}} What are the common \emph{design methods and approaches} employed in smart building research to shape human experiences? 
    \item \textbf{\textit{RQ3.}} What types of \emph{technological interventions} are developed or conceptualised to enhance human interactive experiences in buildings?
\end{itemize}


\subsection{Framework Stage 2: Identifying Relevant Studies}
To ensure that our scoping review investigates disciplines beyond the HCI domain, we queried for works in two established databases. The ACM Digital Library (DL) primarily features computer science and engineering research (including HCI). Scopus offers multidisciplinary coverage, with frequent indexing of architecture works, making it a valuable resource for architecture studies on smart building research\sidenote{List of architecture-based sources available here when filtered by subject area: \url{https://www.scopus.com/sources.uri}}. 

\paragraph{Pilot Searches}
We first conducted pilot searches on \textit{smart buildings} and \textit{experiences}, which resulted in some 9000 references in Scopus alone. We reviewed the first 20\%, finding many irrelevant studies that did not involve human experiences, such as Internet of Things (IoT) and civil and computational engineering. Based on this, we decided to include \textit{interaction} (\texttt{interac*}) in our search strategy to refer to interactive human experiences in smart buildings.  



\paragraph{Employed Search Queries}
Our search queries are explained in \autoref{fig:queries}. The search involved seven queries (Query~1--Query~7), structured around three keyword categories: 1) environment/domain-specific descriptors, 2) experience, and 3) interaction keywords. We conducted the searches using a combination of the keyword groups. A combination with ``built environment'' was searched separately due to the large volume of results, while more specific environmental descriptors---(b) in \autoref{fig:queries}---were combined in a second search. Domain-specific descriptors were handled individually, and when these already included interaction (e.g. ``human-building interaction''), only experience keywords were added. An example of a search query (applied to Scopus) is: \texttt{"built environment*" AND experien* AND interac*}.  The complete list of queries we used are included in the supplementary materials. 

We conducted the searches based on papers' keywords, abstract, and title across both databases\sidenoteup{34pt}{Scopus supports a combined search of titles, abstracts, and keywords. In contrast, the ACM DL requires these parameters to be searched separately. To ensure consistency and comparability, we conducted separate searches within the ACM DL: first by titles, then by abstracts, and finally by keywords, before collating the results.}. This generated \textbf{5,526 references} across both sources, collected during April -- May 2024 (later increased by 22 papers through snowballing \cite{badampudi2015experiences}, see \autoref{fig:flowdiagram}).


\begin{widefigure}[tbp]
    \includesvg[width=\linewidth]{images/methodology/methodology-search.svg}
    \caption{The search comprised seven queries (Query 1 -- Query 7), combining nine components (see (a)--(i)). The nine components consisted of either 1) environment/domain-specific descriptors (first column), 2) a keyword for experience (second column), or 3) a keyword for interaction (third column). We conducted the searches using combinations of these---query list on the right. Thus, environmental descriptors were combined with experience and interaction keywords (red line). This was implemented for "Built environment" separately---(a) in red box---as this yielded many results, while more specific terms were combined in a second search---(b) in red box. Domain-specific descriptors were handled separately again---(c) and (d) in green box; green line. When the domain-specific descriptor already included the word interaction (e.g. 'human-building interaction'), the query was only accompanied by the experience keyword---(e), (f), and (g) in blue box.}
    \label{fig:queries}
\end{widefigure}

\subsection{Framework Stage 3: Study Selection} 
\label{sec-ch2:iecrit}
During our pilot search, we developed initial inclusion-exclusion (I/E) criteria for the first (title and abstract) screening stage, refining them as we grew familiar with the literature. Broadly, references were included if they: (1) were set within the context of a smart building, (2) discussed human experiences, and (3) explored the use or potential use of technologies. References solely focused on engineering, algorithms, or building design without addressing human interaction or the use of technology were excluded. In addition, we excluded papers that were not in English and secondary research (e.g. meta-reviews). The first author applied the I/E criteria to the 5,548 citations, reading all associated abstracts, keywords, and titles. If the relevance of a study was unclear from the abstract, the study was retained for a full article review. This process resulted in 368 shortlisted abstracts. After a review by the last author and the removal of duplicate studies, 338 abstracts were shortlisted for full-paper screening. For the latter, the criteria established for inclusion and exclusion were specified and expanded:

\begin{enumerate}
    \item We included studies around \emph{human interaction} within smart buildings using \emph{any form of technology} for occupants. Conversely, we excluded studies that focused on technology development without clear relevance to user interaction or experience, such as works on theoretical models, architectural frameworks, or technical optimisations. 
    \item Our research focuses on \emph{real-world environments}, however, we incorporated papers that utilise Virtual Reality (VR), Augmented Reality (AR), and Mixed Reality (MR) to augment existing spaces, enhancing spatial experiences and smart building capacities~\cite{albarrak2023using, knierim2019exploring}. However, we excluded research focused on designing purely virtual environments or those centred on gaming realism and virtual worlds without connection to physical space.
    \item We included papers that examined human experiences \emph{within smart buildings and building fa\c{c}ades}. However, we excluded papers focused on outdoor experiences like hiking, cycling, and city walks. Additionally, automotive contexts were excluded as they fall outside the scope of smart buildings.
    \item We excluded papers on domotics, mobile apps, and Conversational User Interfaces (CUIs) that lacked a focus on spatial factors of human experience or interaction within a smart building. Studies on human-robot interaction without connection to physical space were also excluded.
\end{enumerate}

The first and last authors held discussions to resolve any points of confusion throughout the study selection. After reading the articles in full and cross-checking, references cited by the selected papers that may have been accidentally left out, a total of \textbf{192 articles} were selected for inclusion in the scoping review. 

\subsection{Framework Stage 4: Charting the data}
The next stage involved charting key information items from the selected papers. Our charting approach was akin to a ``descriptive-analytical narrative review'' \cite{arksey2005scoping}, where a common analytical framework was applied to the papers, enabling us to systematically collect and analyse standard information from each study. The first and last authors collaboratively developed a charting table template to capture information related to the research questions. This process involved categorisation (labelling) and extraction of relevant sentences for later analysis. While some papers were straightforward to chart, others involved multiple discussions to satisfactorily capture the research perspectives. 


Our template recorded \textit{title}, \textit{author(s)}, \textit{publication year}, \textit{publication source/venue}, \textit{publication type} and \textit{type of smart building} addressed. We structured our template to record items directly related to our research questions. We documented \textit{research motivations}, and \textit{methodologies} including research and design methodologies to infer whether the research was qualitative, quantitative, or mixed-methods. We recorded the \textit{interactive technologies} and any \textit{instruments} used to assess the environments or their users. Central to our research, we noted the \textit{aspects of human experience} discussed. \textit{Key findings} and the papers' \textit{visions} for future were also recorded (in case we needed to examine evidence supporting claims made by the studies), but were considered out of scope for this paper. The charting template with an example entry is included in Appendix Section \ref{sec:appendix-scr-template}. Throughout this iterative process, the research team met regularly to discuss the scope of the review, refine definitions, and ensure consistency in the categorisation of information.


\subsection{Framework Stage 5: Collating, Summarising, and Reporting}
After charting the study data, we analysed it in two main stages: a descriptive summary followed by an interpretive analysis (semantic and latent approaches) \cite{byrne2021worked}. First, we examined the characteristics of the literature such as the venues they were published at. Next, we adopted a hybrid latent coding approach, based on thematic analysis \cite{braun2023thematic}, to analyse the charting table. This approach was hybrid because while we employed an inductive process to create new codes and themes, we also incorporated deductive coding for the four dimensions of comfort that we adopted from existing literature~\cite{alavi2017comfort} as well as \citet{brand1995buildings}'s shearing layers (details described below). The coding was highly iterative and completed over multiple rounds by the first author, with intermittent discussions with the full research team, eventually resulting in overarching themes for the targeted human experiences. Though targeted human experiences are the resulting themes, the other core categories (design mechanisms and technological interventions) were iterated on as well. Our method borrows from reflexive thematic analysis in the organic, flexible nature of the rounds of iteration of codes and themes, but overall is better aligned with a codebook approach~\cite{braun2022conceptual}.

In line with our research questions, we identified key themes around human experiences in smart buildings by coding papers based on the \emph{primary} experiences they targeted. After identifying the targeted experiences, we recognised that the papers could also be analysed for the design measure or approach (hereafter referred to as design mechanism) used to achieve that experience. We therefore developed RQ2 and conducted a second round of this analysis where we coded each paper for the \emph{primary} design mechanisms employed to achieve the targeted human experiences. To address our third research question, we categorised the technology data by type, iteratively grouping them based on observed commonalities and patterns. We also coded the papers based on their technologies' positioning in \citet{brand1995buildings}'s shearing layers which describe distinct layers of a building based on its temporal lifespan \cite{brand1995buildings}.  Finally, the full research team discussed and finalised the resulting themes of targeted human experiences, design mechanisms, and types of technological interventions over two meetings. These discussions considered the suitability and terminology of these themes in the form of the example sentences presented in \autoref{fig:pieces}. The full list of corpus papers and the analysis process are documented in the supplementary materials.

\pgfdeclarelayer{background}
\pgfsetlayers{background,main}

%colorbrewer 3-class RdYlBu
\definecolor{thisred}{RGB}{252,141,89} 
\definecolor{thisyellow}{RGB}{255,255,191} 
\definecolor{thisblue}{RGB}{145,191,219} 



\begin{figure*}[t]
{
\footnotesize
\resizebox{\textwidth}{!}{
\begin{tikzpicture}
    \tikzstyle{block}=[
        draw,
        node distance=0.5cm,
        rounded corners,
        thick,
        text width= 2cm,
        minimum height=.8cm,
        align=center,
        font=\small,
        every node/.style={inner sep=0,outer sep=0},
    ]

    \tikzstyle{header}=[rotate=90, align=center]
    \tikzstyle{line}=[draw, thick]%-latex for arrow lines


    \node[] (statement1) {\emph{In this paper,}};
    \node[block, right = .1cm of statement1, fill = thisred] (techinterv) {technological intervention};

    \node[right = .3cm of techinterv] (statement2pre) { };
    \node[above =.01cm of statement2pre] (statement2preabove) {\emph{is}};
    \node[below =.01cm of statement2pre] (statement2prebelow) {\emph{are}};

    \node[right = .3cm of statement2pre] (verbs1) {\emph{explored}}; 
    \node[above = .05cm of verbs1] (verbs1above) {\emph{used}}; 
    \node[below = .05cm of verbs1] (verbs1below) {\emph{envisioned}};
    \node[below = .05cm of verbs1below] (ellips1) {\emph{...}};
    \node[right = .1cm of verbs1] (statement2post) {\emph{as a way to}}; 

    \node[right = .1cm of statement2post] (verbs2) {\emph{support}}; 
    \node[above = .05cm of verbs2] (verbs2above) {\emph{improve}}; 
    \node[below = .05cm of verbs2] (verbs2below) {\emph{impact}};
    \node[below = .05cm of verbs2below] (ellips2) {\emph{...}};

    \node[right = .1cm of verbs2] (statement2postpost) {\emph{users'}};

    \node[block, right = .1cm of verbs2, fill = thisblue] (the) {targeted human experience};

    \node[right = .1cm of the] (statement3) {\emph{through}}; 

    \node[block, right = .1cm of statement3, fill = thisyellow] (designmech) {primary design mechanism}; 
    
    \node[right = .1cm of designmech] (colon) {\emph{.}}; 

    %\node[below = of techinterv] (statement1) {\emph In this paper, ... is used explored envisioned as a way to improve support influence users' ...  through ...};

    % Draw edges
    %\path [line] (RQs) -- (protocol);
    %\path [line] (protocol) -- (prereg);
    %\path [line] (prereg) -- (search);
    %\path [line] (search) -- (inclExcl);
    %\path [line] (inclExcl) -- (appraisal);
    
    % \path [line] (linecoding) -- (themes);
    % \path [line] ([xshift=.3cm]linecoding.south west) -- ([xshift=.3cm, yshift=-.92cm]linecoding.south west) -- (1storder.west);
    % \path [line] ([xshift=.3cm, yshift=-.92cm]linecoding.south west) -- ([xshift=.3cm, yshift=-2.25cm]linecoding.south west) -- (initialrealism.west);
    % \path [line] ([xshift=-1cm]linecoding.south east) -- ([xshift=-1cm, yshift=-.92cm]linecoding.south east) -- (2ndorder.west);
    
    % \path [line] (themes) -- (deductive);
    % \path [line] (deductive) -- (mapping);
    % %\path [line] (themes) -- (mapping);
    % \path [line] ([xshift=-.7cm]themes.south) -- ([xshift=-.7cm, yshift=-.92cm]themes.south) -- (3rdorder.west);
    
    % \path [line] (theory1) -- (deductive);
    % \path [line] (theory2) -- (deductive);
    % \path [line] (initialrealism.east) -- (mapping.west);
\end{tikzpicture}
}}
\caption{To examine the suitability and the terminology of the three core findings---targeted human experiences, primary design mechanisms, and technological interventions---we structured them in a sentence, illustrated in this diagram. For example, a paper might be described in the form of ``In this paper, \textbf{Environmental Sensing} is used as a way to support \textbf{Comfort} through \textbf{Personalisation}''.}\label{fig:pieces}
\Description{...}
\end{figure*}

